% Options for packages loaded elsewhere
\PassOptionsToPackage{unicode}{hyperref}
\PassOptionsToPackage{hyphens}{url}
\documentclass[
 12pt,
]{ctexart}
\usepackage{xcolor}
\usepackage{amsmath,amssymb}
\setcounter{secnumdepth}{-\maxdimen} % remove section numbering
\usepackage{iftex}
\ifPDFTeX
 \usepackage[T1]{fontenc}
 \usepackage[utf8]{inputenc}
 \usepackage{textcomp} % provide euro and other symbols
\else % if luatex or xetex
 \usepackage{unicode-math} % this also loads fontspec
 \defaultfontfeatures{Scale=MatchLowercase}
 \defaultfontfeatures[\rmfamily]{Ligatures=TeX,Scale=1}
\fi
\usepackage{lmodern}
\ifPDFTeX\else
 % xetex/luatex font selection
\fi
% Use upquote if available, for straight quotes in verbatim environments
\IfFileExists{upquote.sty}{\usepackage{upquote}}{}
\IfFileExists{microtype.sty}{% use microtype if available
 \usepackage[]{microtype}
 \UseMicrotypeSet[protrusion]{basicmath} % disable protrusion for tt fonts
}{}
\makeatletter
\@ifundefined{KOMAClassName}{% if non-KOMA class
 \IfFileExists{parskip.sty}{%
 \usepackage{parskip}
 }{% else
 \setlength{\parindent}{0pt}
 \setlength{\parskip}{6pt plus 2pt minus 1pt}}
}{% if KOMA class
 \KOMAoptions{parskip=half}}
\makeatother
\usepackage{longtable,booktabs,array}
\usepackage{calc} % for calculating minipage widths
% Correct order of tables after \paragraph or \subparagraph
\usepackage{etoolbox}
\makeatletter
\patchcmd\longtable{\par}{\if@noskipsec\mbox{}\fi\par}{}{}
\makeatother
% Allow footnotes in longtable head/foot
\IfFileExists{footnotehyper.sty}{\usepackage{footnotehyper}}{\usepackage{footnote}}
\makesavenoteenv{longtable}
\setlength{\emergencystretch}{3em} % prevent overfull lines
\providecommand{\tightlist}{%
 \setlength{\itemsep}{0pt}\setlength{\parskip}{0pt}}
\usepackage{bookmark}
\IfFileExists{xurl.sty}{\usepackage{xurl}}{} % add URL line breaks if available
\urlstyle{same}
\hypersetup{
 hidelinks,
 pdfcreator={LaTeX via pandoc}}

\author{SCX}
\date{2026}

\begin{document}

\section{SCX Framework: The Next Theorems(The Next
Theorems)}\label{scxux6846ux67b6ux4e0bux4e00ux4ee3ux5b9aux7406ux7fa4the-next-theorems}

\subsection{Draft by Theoretical Architect \textbar{}
June 29, 2026}\label{ux7406ux8bbaux67b6ux6784ux5e08ux8349ux6848-2026ux5e746ux670829ux65e5}

\begin{center}\rule{0.5\linewidth}{0.5pt}\end{center}

\begin{quote}
\textbf{Premise}:T1-T7all hold.T1-T4(static audit completeness),T5(active learning bounds),T6(protocol game equilibrium),T7(cross-domain transfer fidelity),andSpring
SE-1 (dynamic evolution) and Situs (spatial encoding) form a completefoundational theoretical system.
\end{quote}

\begin{center}\rule{0.5\linewidth}{0.5pt}\end{center}

\subsection{一, because果发现Theorem:labelnoise
vs.~because果混淆can分离性}\label{ux4e00ux56e0ux679cux53d1ux73b0ux5b9aux7406ux6807ux7b7eux566aux58f0-vs.-ux56e0ux679cux6df7ux6dc6ux7684ux53efux5206ux79bbux6027}

\subsubsection{(Causal Discovery Theorem ---
CD-Theorem)}\label{causal-discovery-theorem-cd-theorem}

\subsubsection{Context}\label{ux8bedux5883}

Theorem
Theorem 3 (Uncertainty Principle) asserts: within the SCX audit framework,\textbf{noise and intrinsic difficulty are indistinguishable}.Formally, given the Cercis
Score S = Q +
ηN,审计者cannot beonlyFromS观测值分解出the quality componentQandnoise分量N------exist in一个不can约information壁垒.

butTheorem
3left a key gap unclosed:\textbf{noise(labelnoise, 测量noise)andnot yet观测混杂because子(unobserved
confounders)whether也不can区分?}
Thisis一个level更深Problem.labelnoiseisauditing监督signalrandom扰动;because果混淆isbecausenot yet观测共同原because导致虚假关联.两者Under观测distributiononcan能表现出completely相同statisticalfeature,butit们干预Meaning截然不同------before者canviaData清洗缓解,after者mustviabecause果identify(do-calculus, toolvariableetc.)Resolution.

\subsubsection{Theorem陈述(Draft)}\label{ux5b9aux7406ux9648ux8ff0ux8349ux6848}

\begin{quote}
\textbf{CD-Theorem(because果can分离Theorem)}: Let D = \{(xᵢ, yᵢ)\}
provides审计Data集,yᵢ = f*(xᵢ) + εᵢ + β·zᵢ,where εᵢ providesIndependentlabelnoise,zᵢ
providesnot yet观测混杂because子,β provides混杂强度toward量.definition审计者canacquireinformation集provides I\_A =
\{S(x), ∇S(x), H\_S(x)\}(Cercis Score及its一阶, 二阶gradient).

\textbf{(i) Weak Separability(Open Problem)}:onlymakinguses I\_A,exist in一个检验statistical量
τ(I\_A) makingwe getas β ≠ 0 ,τ
势function(power)严格大于only含labelnoise情形under显著性水平,i.e.,: P(τ
\textgreater{} c\_α \textbar{} β ≠ 0) \textgreater{} α for some β ≠ 0 and
ε bounded. theConclusionwhetherprovides真,depending on η(Sprintnoise parameter)whetherand z
 I(η; z) \textgreater{} 0.

\textbf{(ii) Strong Separability(Rigorously derivable,FromTheorem 3直接导出)}:If η and z
Independent(i.e.,I(η; z) =
0),thenlabelnoiseandbecause果混淆UnderSCXframeworkwithin\textbf{absolutely indistinguishable}.审计者cannot be构建任何优于random猜测classify器来区分
ε-dominated samplesand z-dominated samples.

\textbf{(iii) 继Theorem(Rigorously derivable,FromT7导出)}:Ifexist in一个源域 D\_s
satisfybefore门准then(front-door criterion),and T7 跨域迁移
ε-保真度conditionholds,thenUndertarget域 D\_t ,the confounding effect β Cercis bias ‖ΔS‖
we haveUpper bound: ‖ΔS‖ ≤ g(ε, δ) + O(1/√n) where δ provides源域andtarget域distributionDifference Situs
distance between time slices.
\end{quote}

\subsubsection{Annotation}\label{ux6807ux6ce8}

\begin{longtable}[]{@{}ll@{}}
\toprule\noalign{}
Part & Status \\
\midrule\noalign{}
\endhead
\bottomrule\noalign{}
\endlastfoot
(i) Weak Separability & \textbf{Open Problem} --- depending on η-z depended onstructure实证发现 \\
(ii) Strong Separability & \textbf{Rigorously derivable} --- Theorem 3 direct corollary \\
(iii) 继Theorem & \textbf{Rigorously derivable} --- T7 + Situsdistance between time slices三角不etc.式 \\
\end{longtable}

\subsubsection{coreintuition}\label{ux6838ux5fc3ux76f4ux89c9}

noiseis''Nonestructuredifficult'',混杂is''we havestructuredifficult''.Theorem
3说UnderSlevel两者不can分,but\textbf{SHigh阶structure(gradient场∇Scurl, Hessian秩)can能encoding了because果information}------Thisis(i)开放Partcore赌注.IfSitusencoding保留了because果图DAG骨架information,then我们就we have楔子撬开This区分.

\begin{center}\rule{0.5\linewidth}{0.5pt}\end{center}

\subsection{二, 联邦审计Theorem:Zero-Knowledge Multi-Party Audit Completeness}\label{ux4e8cux8054ux90a6ux5ba1ux8ba1ux5b9aux7406ux96f6ux77e5ux8bc6ux591aux65b9ux5ba1ux8ba1ux5b8cux5907ux6027}

\subsubsection{(Federated Audit Theorem ---
FA-Theorem)}\label{federated-audit-theorem-fa-theorem}

\subsubsection{Context}\label{ux8bedux5883-1}

``Audit Sword''(Audit
Sword)确立了SCXcore权力:\textbf{任何SCXmakinguses行providescanis bounded below byIndependent审计}.butThis剑目beforeAssume审计者拥we haveCompleteData访Question权.Under联邦scenariounder------多方持we have私密Data, onlyviagradientor摘要statistical交互------Audit Sword遇到了关节锁:你such as何审计一个你cannot be看seeData集?

Federated learning already has tools like安全聚合(Secure
Aggregation), 差分隐私(DP), 多方compute(MPC).butit们Resolutionis''防止泄露'',while非''makingcan审计''.SCXtargetis逆toward:我们要Under\textbf{不暴露The originalDatabefore提under,making审计Conclusion具we haveand全Data审计相同statistical效力}.

\subsubsection{Theorem陈述(Draft)}\label{ux5b9aux7406ux9648ux8ff0ux8349ux6848-1}

\begin{quote}
\textbf{FA-Theorem(联邦审计Theorem)}: Suppose there are K
data parties, each holding a private dataset D\_k, k = 1,\ldots,K.联邦聚合器viaprotocol Π
Output全局Model M\_Π.审计者 A canacquireinformation集provides: I\_Fed = \{S\_k, ∇S\_k,
H\_\{S,k\}\}\_\{k=1\}\^{}K ∪ \{M\_Π\} where S\_k provideseach方 Cercis
Score(computable locally and shareable with encryption).

\textbf{(i) Federated Audit Completeness(Open Problem,Partcan推)}:
Define the full-information audit statistic T\_full = T(D\_1,\ldots,D\_K),and the federated audit statistic
T\_fed = T\_F(I\_Fed).Ifprotocol Π satisfy: (a) each方提交 S\_k
is verified for honesty via a zero-knowledge proof(i.e., ∃ π\_ZK : Verify(π\_ZK, S\_k, Commit(D\_k))
= 1); (b) SitusencodingUndereach方间maintain同态:Situs(∪D\_k) ≅ ⊕\_k
Situs(D\_k)(⊕providessome种can交换组合operator), thenexist in一个 ε-audit equivalence:
\textbar Power(T\_fed) - Power(T\_full)\textbar{} ≤ ε(K, δ) where δ
provideseach方差分隐私预算(Ifnot yetmakingusesDP,δ = 0 and ε → 0 as n → ∞).

\textbf{(ii) informationlower bounded(Rigorously derivable,FromTheorem 3导出)}: If ∃ k makingwe get D\_k
labelnoise parameter η\_k
is completely unknown to the auditor(i.e.,Noneπ\_ZKverification),thenauditing任何and the federated audit statistic
T\_fed,exist in一个混淆策略makingwe get: sup\_\{D\} \textbar E{[}T\_fed{]} -
E{[}T\_full{]}\textbar{} ≥ const · η\_k ·
\textbar D\_k\textbar/\textbar D\_total\textbar{}
i.e.,审计效力Lossandnot yetverificationnoiseUnder总Data占比成正比.

\textbf{(iii) Telegraph Public Goods Extension(Open Problem)}:
If''电报公共品''(Datastoreisfundamental设施)延伸至联邦层,i.e.,exist in一个去心化store层承诺(commitment)了allwe have
D\_k Situs hash,thenFA-Theorem (i) π\_ZK
要求canis bounded below byreplaceprovidesstoreproof(Proof of Storage),ε Upper bound进一步收紧.
\end{quote}

\subsubsection{Annotation}\label{ux6807ux6ce8-1}

\begin{longtable}[]{@{}
 >{\raggedright\arraybackslash}p{(\linewidth - 2\tabcolsep) * \real{0.5000}}
 >{\raggedright\arraybackslash}p{(\linewidth - 2\tabcolsep) * \real{0.5000}}@{}}
\toprule\noalign{}
\begin{minipage}[b]{\linewidth}\raggedright
Part
\end{minipage} & \begin{minipage}[b]{\linewidth}\raggedright
Status
\end{minipage} \\
\midrule\noalign{}
\endhead
\bottomrule\noalign{}
\endlastfoot
(i) Federated Audit Completeness & \textbf{Mixed} ---
Situs同态组合Partis开放(requires构造具体组合operator);ε界DerivationGivenZKconditionunderis严格 \\
(ii) informationlower bounded & \textbf{Rigorously derivable} --- Theorem 3 +
information-theoretic data processing inequality \\
(iii) Telegraph Public Goods Extension & \textbf{Open Problem} ---
depending onstore层design,理论onwe have清晰路径 \\
\end{longtable}

\subsubsection{coreintuition}\label{ux6838ux5fc3ux76f4ux89c9-1}

联邦审计essential:\textbf{安全聚合保护了Data方免受审计者窥探,butAudit Sword要求审计者能看穿聚合}.Resolution路径不Under密码学level(那会导致''全we haveor全None''),whileUnderSCXinformation架构level------Situsencodingasprovides''审计指纹'',Under保留can审计性同丢弃了can重构性.This就isZKproof介入楔入点.

\begin{center}\rule{0.5\linewidth}{0.5pt}\end{center}

\subsection{三, auditing抗RobustnessTheorem:auditing抗扰动is''difficult''Special case还is第三类?}\label{ux4e09ux5bf9ux6297ux9c81ux68d2ux6027ux5b9aux7406ux5bf9ux6297ux6270ux52a8ux662fux56f0ux96beux7684ux7279ux4f8bux8fd8ux662fux7b2cux4e09ux7c7b}

\subsubsection{(Adversarial Robustness Theorem ---
AR-Theorem)}\label{adversarial-robustness-theorem-ar-theorem}

\subsubsection{Context}\label{ux8bedux5883-2}

Theorem 3建立了noiseanddifficult二元不can区分framework.butauditing抗sample(adversarial
examples)发现toThis二元framework制造了麻烦:auditing抗扰动不israndomnoise------itis针auditingModelgradient定toward最optimizeResult;it也不is''自然difficult''------humans can effortlessly classify adversarial examples correctly.This raises a fundamental question:

\textbf{SCXCercisclassify学,whetherexist in''auditing抗性''asprovides第三basicclass,andQ(质量), N(noise)alongside,i.e.,
S = Q + ηN + γA?}

\subsubsection{Theorem陈述(Draft)}\label{ux5b9aux7406ux9648ux8ff0ux8349ux6848-2}

\begin{quote}
\textbf{AR-Theorem(auditing抗RobustnessTheorem)}: Let the adversarial perturbation δ\_adv =
argmax\_\{‖δ‖≤ε\} L(f(x+δ), y),and ‖δ\_adv‖
UnderSCXaudit frameworkwithin受Situsdistance between time slicesmetric.definitionAdversarial exposure度(Adversarial
Exposure)A(x) = ‖∇\_x S(x)‖ · ε / S(x).

\textbf{(i) 还原Theorem(Open Problem,Partcan推)}: IfModel f Cercisgradient场 ∇S
UnderxSitus邻域withinsatisfyLipschitzconditionandLipschitzconstant L \textless{}
S(x)/ε,thenauditing抗扰动canis bounded below byTheorem 3吸收到''noise''分量: A(x) ≤
η\_eff(x),where η\_eff = η + L·ε/S(x)(effective noise rate). Conversely,If L ≥
S(x)/ε,thenauditing抗扰动构成一个不canis bounded below byTheorem 3吸收\textbf{第三类奇点}.

\textbf{(ii) Audit Detectability(Rigorously derivable,FromT5主动学习Theorem导出)}:
definition审计者auditing抗detectstatistical量: D\_adv(x) = ‖∇S(x)‖ /
Var\_\{x'∈N(x)\}{[}S(x'){]} where N(x) provides Situs local neighborhood.If γ
\textgreater{}
0(i.e.,exist in第三类auditing抗分量),thenas审计者采usesT5主动学习策略以is maximized化I(S;
D\_adv),detectsampleComplexityprovides: n\_detect ≥ Ω(1/γ² · log(1/δ))
andthelower boundedauditing最优主动学习策略is紧.

\textbf{(iii) 相变Theorem(Open Problem)}: exist in一个临界Cercisthreshold
S\_crit,makingwe get: - as S(x) \textgreater{}
S\_crit:auditing抗扰动provides''difficult''Special case(reduced to theTheorem 3framework), - as S(x) ≤
S\_crit:auditing抗扰动表现provides第三类,andA分量andN分量Under∇Sspace正交. S\_crit
解析形式及itsandModel架构relationshipis开放.
\end{quote}

\subsubsection{Annotation}\label{ux6807ux6ce8-2}

\begin{longtable}[]{@{}
 >{\raggedright\arraybackslash}p{(\linewidth - 2\tabcolsep) * \real{0.5000}}
 >{\raggedright\arraybackslash}p{(\linewidth - 2\tabcolsep) * \real{0.5000}}@{}}
\toprule\noalign{}
\begin{minipage}[b]{\linewidth}\raggedright
Part
\end{minipage} & \begin{minipage}[b]{\linewidth}\raggedright
Status
\end{minipage} \\
\midrule\noalign{}
\endhead
\bottomrule\noalign{}
\endlastfoot
(i) 还原Theorem & \textbf{Open Problem} ---
Lipschitzconditionis充分butcan能非必要;γ\textgreater0区域严格feature化待complete \\
(ii) Audit Detectability & \textbf{Rigorously derivable} --- T5主动学习framework直接shoulduses \\
(iii) 相变Theorem & \textbf{Open Problem} ---
相变exist in性proofandS\_crit解析形式allnot yetResolution \\
\end{longtable}

\subsubsection{coreintuition}\label{ux6838ux5fc3ux76f4ux89c9-2}

auditing抗扰动UnderCercisspace留under指纹不isS绝auditing值(becauseprovidesitcan能很小),whileis\textbf{∇SUnderlocal neighborhoodwithin异常structure}------正常difficultsample∇Svaries smoothly,whileauditing抗sample周围exist ingradient断裂.This suggests that theA分量can能Under∇S拓扑不variable(such as局部曲率, Hessian谱)encoding,while非UnderSitself.

\begin{center}\rule{0.5\linewidth}{0.5pt}\end{center}

\subsection{四, 序SCXTheorem:非平稳distributionunder审计drift}\label{ux56dbux65f6ux5e8fscxux5b9aux7406ux975eux5e73ux7a33ux5206ux5e03ux4e0bux7684ux5ba1ux8ba1ux6f02ux79fb}

\subsubsection{(Temporal SCX Theorem ---
TS-Theorem)}\label{temporal-scx-theorem-ts-theorem}

\subsubsection{Context}\label{ux8bedux5883-3}

staticSCXframeworkAssumeDatadistributionP(X,Y)固定.Spring SE-1引入了dynamic evolution(model update →
Datadistributiondrift →
SCXparameterηdrift),butSE-1目before只process了''model updateasprovidesoutside生冲击''scenario.现实distributiondrift更复杂:
- \textbf{概念drift}(concept drift):P(Y\textbar X)changes; -
\textbf{协variabledrift}(covariate shift):P(X)changes; -
\textbf{priordrift}(label shift):P(Y)changes.

更criticalis,\textbf{Sprintcoreparameterη(noise-difficulty coupling coefficient)自身as间drift}------This就is''审计审计''Problem:谁Under审计Sprintitself?

\subsubsection{Theorem陈述(Draft)}\label{ux5b9aux7406ux9648ux8ff0ux8349ux6848-3}

\begin{quote}
\textbf{TS-Theorem(序SCXTheorem)}: Let间索引 t =
0,1,\ldots,T.Undereach个刻 t,define the data distribution P\_t(X,Y),auditingshouldCercis Score
S\_t,Sprintparameter η\_t,andSitusencoding U\_t = Situs(P\_t).Assume η\_t
evolution遵循: η\_\{t+1\} = η\_t + Δη(P\_t, M\_t) + ξ\_t where Δη
providesdeterministic性drift(guaranteed byModel M\_t andDatadistribution P\_t 共同decide to),ξ\_t
providesrandom冲击(mean 0, variance σ²\_ξ).

\textbf{(i) driftcandetect性(Rigorously derivable,FromT7跨域迁移导出)}:
If两刻Situsdistance between time slices d\_Situs(U\_t, U\_\{t+1\}) ≥ δ\_min,thenUnder显著性水平
α the auditor can detect ≥ 1-β drift in η drift,allrequiressample量satisfy: n\_t +
n\_\{t+1\} ≥ 2(σ²\_ξ + κ·δ²\_min)·(z\_\{α/2\} + z\_β)² / δ²\_min where κ
providesSitusdistance between time slicesandηdriftLipschitzconstant(FromT7继承).

\textbf{(ii) drift归because(Open Problem)}: asdrift in η driftafter,能否willits分解provides:
Δη = Δη\_concept + Δη\_covariate + Δη\_label + Δη\_model
i.e.,归because至概念drift, 协variabledrift, priordriftandmodel updateeach自contribution?IfNone额outsideAssume(such asPartcan观测verification集),the分解UnderSCXframeworkwithinis\textbf{不canidentify}.然while,Ifexist in一组anchorssample
A ⊂ X satisfy P\_t(Y\textbar X∈A) ≈
P\_\{t+1\}(Y\textbar X∈A)(概念stableanchors),then Δη\_concept canis bounded below byIndependentestimate.

\textbf{(iii) Sprintself-audit limitsTheorem(Open Problem,Partcan推)}: Sprint
SE-1自我CorrectedMechanism(η → η'
viaModel重Training)whether总能willηback to equilibrium?Define a Lyapunov function V(η) = (η -
η\emph{)\^{}2,where η} provides最优noise parameter.Ifexist in一个紧集 Ω ⊂ {[}0,1{]}
makingwe getauditing η\_t ∈ Ω,we have: E{[}V(η\_\{t+1\}) \textbar{} η\_t{]} ≤
ρ·V(η\_t),ρ \textless{} 1 thenSprint过程is几何遍历(geometrically
ergodic),η\_t converges to η*.butIfdistributiondriftrate v\_drift satisfy v\_drift
\textgreater{}
(1-ρ)/Δt,thenSprintdiverges------自审计能力we have\textbf{rateon限}.
\end{quote}

\subsubsection{Annotation}\label{ux6807ux6ce8-3}

\begin{longtable}[]{@{}
 >{\raggedright\arraybackslash}p{(\linewidth - 2\tabcolsep) * \real{0.5000}}
 >{\raggedright\arraybackslash}p{(\linewidth - 2\tabcolsep) * \real{0.5000}}@{}}
\toprule\noalign{}
\begin{minipage}[b]{\linewidth}\raggedright
Part
\end{minipage} & \begin{minipage}[b]{\linewidth}\raggedright
Status
\end{minipage} \\
\midrule\noalign{}
\endhead
\bottomrule\noalign{}
\endlastfoot
(i) driftcandetect性 & \textbf{Rigorously derivable} --- T7 + 标准Assume检验framework \\
(ii) drift归because & \textbf{Open Problem} ---
NoneanchorsAssume不canidentify(proved);we haveanchors方towardcan行buttechniqueonnot yetcomplete \\
(iii) Sprintself-audit limits & \textbf{Open Problem} ---
Lyapunovconditionis充分but非必要;driftrateUpper bound紧性待proof \\
\end{longtable}

\subsubsection{coreintuition}\label{ux6838ux5fc3ux76f4ux89c9-3}

序SCXessential张力:\textbf{Sprintrequiresηto calibrate the audit,butηitselfUnderis bounded below by审计distribution}.This构成了一个recursive审计Problem------``Audit Sword''Under间Dimensiononrequires自我指涉(self-referentiality).TS-Theoremattempts to formalize:This种自我指涉Under什么conditionunderis良性(convergent),Under什么conditionunderis恶性(diverges).

\begin{center}\rule{0.5\linewidth}{0.5pt}\end{center}

\subsection{五, 人机协同审计Theorem:Theorem
3认知boundary}\label{ux4e94ux4ebaux673aux534fux540cux5ba1ux8ba1ux5b9aux7406theorem-3ux7684ux8ba4ux77e5ux8fb9ux754c}

\subsubsection{(Human-AI Collaborative Audit Theorem ---
HC-Theorem)}\label{human-ai-collaborative-audit-theorem-hc-theorem}

\subsubsection{Context}\label{ux8bedux5883-4}

Theorem
3Core Claimis:noiseanddifficultUnderSCX自动audit frameworkwithin不can区分.butThis句话里最criticalThe two字is:\textbf{``frameworkwithin''}.audit framework操asboundarydepending on我们to审计者配备了什么样informationprocess器.If审计者只is一个that the optimal selection based onCercis
Score自动化system,thenTheorem
3isiron law.butIf审计者can以is\textbf{人类专家 + AIsystemMixed体}呢?

人类拥we have些什么AI审计者没we have? -
\textbf{Prior causal knowledge}(regarding世界such as何运as非statistical知识); -
\textbf{Counterfactual reasoning ability}(``IfAnnotation者更仔细,Thissamplelabel会变吗?'');
- \textbf{Metacognitive judgment}(``Thissample我都不deterministic,all以itcan能真很难'').

Theorem 3boundary,就is人类认知can以介入并打破noise-difficultauditing称性那条线.

\subsubsection{Theorem陈述(Draft)}\label{ux5b9aux7406ux9648ux8ff0ux8349ux6848-4}

\begin{quote}
\textbf{HC-Theorem(人机协同审计Theorem)}: definition审计systemprovides A = (AI, H,
Π),where AI providesthat the optimal selection based onSCX自动审计module,H provides人类审计者,Π
provides人机交互protocol.Define the human audit budget
B\_H(人类can审查is maximizedsample数),and人类判断function J\_H: X → \{noise,
hard, uncertain\}.

\textbf{(i) 协同增益Theorem(Open Problem,方toward性Conclusion)}: Ifsatisfy以underThe twocondition:
(a) \textbf{information不auditing称condition}:exist in一 samples子集 X\_H ⊂ X,makingwe get人类判断
J\_H(x) andCercis Score S(x) Underconditiondistributiononwe have: P(J\_H=noise \textbar{}
S=s) ≠ P(J\_H=hard \textbar{} S=s) for some s holds
(i.e.,人类判断携带了S之outside额outsideinformation); (b)
\textbf{校准condition}:人类auditing''不deterministic''判断itselfis校准,i.e.,
E{[}J\_H=uncertain \textbar{} S=s{]} as s → 0 ,
then人机协同审计noise-difficult分离效能 Ω\_collab strictly dominates pure AI audit: Ω\_collab
\textgreater{} Ω\_AI = 0(UnderTheorem 3under).

\textbf{(ii) Budget Decay Law(Rigorously derivable,FromTheorem 3导出)}: When the human audit budget
B\_H → 0 ,协同审计退化provides纯AI审计,andconvergentrateprovides: Ω\_collab(B\_H) =
Ω\_AI + c·√(B\_H/n)·Δ\_I(S; J\_H) + o(√(B\_H/n)) where Δ\_I
provides人类判断andCercis Score之间information增量(measured in mutual information).

\textbf{(iii) 失败Theorem------Theorem 3绝auditingboundary(Rigorously derivable)}:
i.e.,making引入任意人类审计能力,以under情形undernoise-difficult仍\textbf{absolutely indistinguishable}:
asandonlyas ∀x, I(J\_H(x); ε(x) \textbar{} S(x)) = 0
i.e.,人类判断andlabelnoiseGivenSconditionundercompletely不携带mutual information.

Equivalently,humans breakTheorem
3充分必要conditionis:exist inat least一 samples,itslabelnoisestatisticalfeatureUnder人类认知表征留under了S之outside痕迹.Ifnoiseis''perfect noise''(and任何can观测featureNone关,including人类prior知识),thenTheorem
3屏障auditing人类同样holds.
\end{quote}

\subsubsection{Annotation}\label{ux6807ux6ce8-4}

\begin{longtable}[]{@{}
 >{\raggedright\arraybackslash}p{(\linewidth - 2\tabcolsep) * \real{0.5000}}
 >{\raggedright\arraybackslash}p{(\linewidth - 2\tabcolsep) * \real{0.5000}}@{}}
\toprule\noalign{}
\begin{minipage}[b]{\linewidth}\raggedright
Part
\end{minipage} & \begin{minipage}[b]{\linewidth}\raggedright
Status
\end{minipage} \\
\midrule\noalign{}
\endhead
\bottomrule\noalign{}
\endlastfoot
(i) 协同增益Theorem & \textbf{Open Problem} ---
direction is clear,butrequires人类认知Experiment来deterministiccondition(a)whetherUnder现实holds(很can能isholds);形式化frameworkalreadycomplete \\
(ii) Budget Decay Law & \textbf{Rigorously derivable} --- Theorem 3 + standard sample efficiency analysis \\
(iii) 失败Theorem & \textbf{Rigorously derivable} --- ThisisTheorem
3逻辑闭包:will''frameworkwithin''extension到''任何cancompute+人类can表征informationprocesssystem'' \\
\end{longtable}

\subsubsection{coreintuition}\label{ux6838ux5fc3ux76f4ux89c9-4}

Theorem
3不is宇宙物理定律,itisSCXinformation架构一个\textbf{透镜boundary}------ittells you,viaCercis
ScoreThis透镜看世界,noiseanddifficult会合provides一体.Human intervention might break this boundary,不isbecauseprovides我们更聪明,whileisbecauseprovides我们via\textbf{不同透镜}------常识because果Model, 具身经验, 社会Context------来看待Data.协同审计essentialisMulti-lens融合.

\begin{center}\rule{0.5\linewidth}{0.5pt}\end{center}

\subsection{总结:Theorem群逻辑structureand成熟度地图}\label{ux603bux7ed3ux5b9aux7406ux7fa4ux7684ux903bux8f91ux7ed3ux6784ux4e0eux6210ux719fux5ea6ux5730ux56fe}

\begin{verbatim}
  T1-T4 (static审计完备)
  / | \
  /  |  \
  T5(主动) T6(博弈) T7(迁移)
  |  |  |
  [AR-Theorem] | [TS-Theorem]
  |  |  |
  +---[CD-Theorem]--+
  |
  [FA-Theorem]
  |
  [HC-Theorem]
\end{verbatim}

\subsubsection{Maturity Summary Table}\label{ux6210ux719fux5ea6ux603bux8868}

\begin{longtable}[]{@{}
 >{\raggedright\arraybackslash}p{(\linewidth - 6\tabcolsep) * \real{0.2222}}
 >{\raggedright\arraybackslash}p{(\linewidth - 6\tabcolsep) * \real{0.3704}}
 >{\centering\arraybackslash}p{(\linewidth - 6\tabcolsep) * \real{0.1852}}
 >{\raggedright\arraybackslash}p{(\linewidth - 6\tabcolsep) * \real{0.2222}}@{}}
\toprule\noalign{}
\begin{minipage}[b]{\linewidth}\raggedright
Theorem
\end{minipage} & \begin{minipage}[b]{\linewidth}\raggedright
Core Claim
\end{minipage} & \begin{minipage}[b]{\linewidth}\centering
Rigorously derivablePart占比
\end{minipage} & \begin{minipage}[b]{\linewidth}\raggedright
Key Open Challenge
\end{minipage} \\
\midrule\noalign{}
\endhead
\bottomrule\noalign{}
\endlastfoot
\textbf{CD-Theorem} & labelnoiseandbecause果混淆canwe have限分离 & \textasciitilde40\%
& η-zmutual information structure;Situsencodingwhether保留because果骨架 \\
\textbf{FA-Theorem} & Zero-knowledge multi-party audit is equivalent to full-data audit &
\textasciitilde50\% & Situshomomorphic composition operator;ZKUnder实际network构造开销 \\
\textbf{AR-Theorem} &
auditing抗性is''difficult''子类还is第三类depending onCercisgradientsmooth性 &
\textasciitilde30\% & the phase transition pointS\_crit解析形式;正交性proof \\
\textbf{TS-Theorem} & ηdriftcandetectbut归becauserequiresanchorsAssume;Sprintwe haveself-audit limits
& \textasciitilde60\% & drift归becauseanchors构造;Lyapunovcondition紧性 \\
\textbf{HC-Theorem} & 人类viaS之outside认知透镜打破Theorem
3,butperfect noise仍不can破 & \textasciitilde55\% &
人类认知information增量实证metric;最优人机交互protocoldesign \\
\end{longtable}

\subsubsection{depended onrelationship}\label{ux4f9dux8d56ux5173ux7cfb}

\begin{itemize}
\tightlist
\item
 CD-Theorem (i) Resolution → 直接enhance TS-Theorem (ii) 归because能力
\item
 AR-Theorem 相变Analysis → provides HC-Theorem
 provide''AIcan自动processauditing抗sample''boundary
\item
 FA-Theorem Situshomomorphism质 → isallwe have联邦/distribution式SCXprotocolfundamental
\item
 HC-Theorem (i) 人机协同增益 → toallwe haveTheoremprovide最终''地面on限''
\end{itemize}

\begin{center}\rule{0.5\linewidth}{0.5pt}\end{center}

\emph{Draft版本 v0.1 \textbar{} SCX理论架构师 \textbar{} 待合as者审阅}

\end{document}
